% Options for packages loaded elsewhere
\PassOptionsToPackage{unicode}{hyperref}
\PassOptionsToPackage{hyphens}{url}
%
\documentclass[
]{article}
\usepackage{amsmath,amssymb}
\usepackage{iftex}
\ifPDFTeX
  \usepackage[T1]{fontenc}
  \usepackage[utf8]{inputenc}
  \usepackage{textcomp} % provide euro and other symbols
\else % if luatex or xetex
  \usepackage{unicode-math} % this also loads fontspec
  \defaultfontfeatures{Scale=MatchLowercase}
  \defaultfontfeatures[\rmfamily]{Ligatures=TeX,Scale=1}
\fi
\usepackage{lmodern}
\ifPDFTeX\else
  % xetex/luatex font selection
\fi
% Use upquote if available, for straight quotes in verbatim environments
\IfFileExists{upquote.sty}{\usepackage{upquote}}{}
\IfFileExists{microtype.sty}{% use microtype if available
  \usepackage[]{microtype}
  \UseMicrotypeSet[protrusion]{basicmath} % disable protrusion for tt fonts
}{}
\makeatletter
\@ifundefined{KOMAClassName}{% if non-KOMA class
  \IfFileExists{parskip.sty}{%
    \usepackage{parskip}
  }{% else
    \setlength{\parindent}{0pt}
    \setlength{\parskip}{6pt plus 2pt minus 1pt}}
}{% if KOMA class
  \KOMAoptions{parskip=half}}
\makeatother
\usepackage{xcolor}
\usepackage[margin=1in]{geometry}
\usepackage{graphicx}
\makeatletter
\newsavebox\pandoc@box
\newcommand*\pandocbounded[1]{% scales image to fit in text height/width
  \sbox\pandoc@box{#1}%
  \Gscale@div\@tempa{\textheight}{\dimexpr\ht\pandoc@box+\dp\pandoc@box\relax}%
  \Gscale@div\@tempb{\linewidth}{\wd\pandoc@box}%
  \ifdim\@tempb\p@<\@tempa\p@\let\@tempa\@tempb\fi% select the smaller of both
  \ifdim\@tempa\p@<\p@\scalebox{\@tempa}{\usebox\pandoc@box}%
  \else\usebox{\pandoc@box}%
  \fi%
}
% Set default figure placement to htbp
\def\fps@figure{htbp}
\makeatother
\setlength{\emergencystretch}{3em} % prevent overfull lines
\providecommand{\tightlist}{%
  \setlength{\itemsep}{0pt}\setlength{\parskip}{0pt}}
\setcounter{secnumdepth}{5}
\usepackage{fontspec}
\setmainfont{DejaVu Serif}
\usepackage{booktabs}
\usepackage{longtable}
\usepackage{float}
\floatplacement{figure}{H}
\usepackage{caption}
\captionsetup{font=small}
\usepackage{booktabs}
\usepackage{longtable}
\usepackage{array}
\usepackage{multirow}
\usepackage{wrapfig}
\usepackage{float}
\usepackage{colortbl}
\usepackage{pdflscape}
\usepackage{tabu}
\usepackage{threeparttable}
\usepackage{threeparttablex}
\usepackage[normalem]{ulem}
\usepackage{makecell}
\usepackage{xcolor}
\usepackage{bookmark}
\IfFileExists{xurl.sty}{\usepackage{xurl}}{} % add URL line breaks if available
\urlstyle{same}
\hypersetup{
  pdftitle={Niðurstöður --- Gróðurbreytingar á náttúruverndarsvæðum Íslands},
  hidelinks,
  pdfcreator={LaTeX via pandoc}}

\title{Niðurstöður --- Gróðurbreytingar á náttúruverndarsvæðum Íslands}
\author{}
\date{\vspace{-2.5em}2026-04-10}

\begin{document}
\maketitle

{
\setcounter{tocdepth}{3}
\tableofcontents
}
\begin{center}\rule{0.5\linewidth}{0.5pt}\end{center}

\section{Minnkun tegundafjölda}\label{minnkun-tegundafjuxf6lda}

Tegundafjöldi minnkaði á milli fyrri og seinni mælingar í 183 af 305
reitum (60\%), á meðan 102 reitir (33\%) sýndu aukningu og 20 héldust
stöðugir á milli athugana. Miðgildi breytingarinnar var -2 tegundir á
reit (meðaltal -2.2). Wilcoxon-próf staðfesti að minnkunin er marktækt
frábrugðin núlli (V = 12,113.5, p \textless{} 0,001), sem þýðir að
hnignunin er raunveruleg og kerfisbundin, ekki tilviljanakennd.

\begin{figure}

{\centering \includegraphics{figures/richness-density-1} 

}

\caption{Mynd 1. Dreifing tegundafjölda við fyrri (blátt) og seinni (rautt) mælingu. Lóðrétta línan sýnir miðgildi fyrri mælinganna.}\label{fig:richness-density}
\end{figure}

\begin{center}\rule{0.5\linewidth}{0.5pt}\end{center}

\section{Endurnýjun tegunda
(Bray--Curtis)}\label{endurnuxfdjun-tegunda-braycurtis}

Endurnýjun tegunda mæld með Bray--Curtis-fjarlægð milli fyrri og seinni
mælinga í hverjum reit. Miðgildi Bray--Curtis-gildis var 0.609 (meðaltal
0.585, n = 163 reitir). Þetta þýðir að þegar bæði fjöldi tegunda og
þekja þeirra eru tekin með í reikninginn, hafði um helmingur
tegundasamsetningarinnar breyst á milli fyrstu og síðustu mælingar ---
að meðaltali á hvern reit. Fylgni milli endurnýjunar tegunda og minnkun
tegundafjölda var marktæk (Spearman ρ = -0.323, p \textless{} 0,001), en
fylgnin var ekki sterk. sem bendir til þess að tvenns konar breytingar
séu í gangi: annars vegar nettótap tegunda og hins vegar breyting á
tegundasamsetningu þar nýjar tegundir skipta öðrum út án þess að
heildarfjöldinn breytist endilega.

\begin{figure}

{\centering \includegraphics{figures/bc-gap-1} 

}

\caption{Mynd 2. Breyting á tegundasamsetningu (Bray–Curtis) í tengslum við tíma milli mælinga. Rauð brotalína = línuleg aðlögun (95\% CI); blá lína = loess-sléttun.}\label{fig:bc-gap}
\end{figure}

Breyting á tegundasamsetningu jókst marktækt með auknum tíma milli
mælinga (hallatala = +0.025 BC-einingar/ár, R² = 0.541, p \textless{}
0,001). Loess-kúrfan sýnir að breytingin er hvað hröðust á fyrstu 15
árunum og hægist síðan á eftir um 20 ár. Þetta bendir til þess að reitir
þar sem langur tími líður á milli mælinga séu ekki endilega þeir sem
hafa breyst mest á öllu tímabilinu.

\begin{center}\rule{0.5\linewidth}{0.5pt}\end{center}

\section{Hvaða umhverfisþættir spá fyrir um
breytingar?}\label{hvauxf0a-umhverfisuxfeuxe6ttir-spuxe1-fyrir-um-breytingar}

Þrír umhverfisþættir voru metnir sem mögulegar skýribreytur: hæð yfir
sjávarmáli, jarðvegsdýpi og fjöldi ára milli mælinga. Hrá
Spearman-fylgni sýnir að hæð og tími milli mælinga fylgjast mjög vel að
(ρ ≈ 0,64--0,70 fyrir BC), en þessi tvær breytur eru ekki óháðar: reitir
á hálendinu höfðu almennt lengra bil milli mælinga. Hlutfylgni, þar sem
áhrif annarra breyta eru hafðar fastar sýnir hvernig þættirnir hegða sér
óháð hverjum öðrum.

\begin{table}[!h]
\centering
\caption{\label{tab:predictor-table}Spearman-fylgni og hlutfylgni þriggja umhverfisþátta við $\Delta$S og BC. Hlutfylgni eru reiknuð með öllum þremur þáttum í líkaninu samtímis.}
\centering
\begin{tabular}[t]{lr>{}rrr>{}rr}
\toprule
Þáttur & ΔS rho (hrátt) & ΔS rho (hlutfylgni) & ΔS p & BC rho (hrátt) & BC rho (hlutfylgni) & BC p\\
\midrule
\cellcolor{gray!10}{Hæð yfir sjávarmáli} & \cellcolor{gray!10}{-0.343} & \textbf{\cellcolor{gray!10}{-0.134}} & \cellcolor{gray!10}{0.0000} & \cellcolor{gray!10}{0.643} & \textbf{\cellcolor{gray!10}{0.170}} & \cellcolor{gray!10}{0.0000}\\
Jarðvegsdýpi & -0.110 & \textbf{-0.188} & 0.0556 & -0.017 & \textbf{0.131} & 0.7675\\
\cellcolor{gray!10}{Bil milli mælinga (ár)} & \cellcolor{gray!10}{-0.349} & \textbf{\cellcolor{gray!10}{-0.106}} & \cellcolor{gray!10}{0.0000} & \cellcolor{gray!10}{0.700} & \textbf{\cellcolor{gray!10}{0.378}} & \cellcolor{gray!10}{0.0000}\\
\bottomrule
\end{tabular}
\end{table}

Þrjár niðurstöður skera sig úr: \textbf{Tími milli mælinga} hefur mest
áhrif á aukna endurnýjun tegunda (BC hlutafylgni ρ = 0.378, p
\textless{} 0,001), sem segir sig að vissu leyti sjálft en staðfestir þó
að gróðurbreytingar eru til staðar og aukast með tíma. \textbf{Hæð yfir
sjávarmáli} er marktæk óháð tíma milli mælinga (ΔS hlutafylgni ρ =
-0.134, BC hlutafylgni ρ = 0.17), sem bendir til þess að hálendisreitir
séu viðkvæmari. \textbf{Jarðvegsdýpi} kemur fram sem sterkasta
sjálfstæða skýribreytan fyrir fyrir \emph{tegundatap} (ΔS hlutafylgni ρ
= -0.188, p \textless{} 0,05) þegar hæð og tími hafa verið festar -
minni jarðvegur tapar fleiri tegundum en djúpur jarðvegur.

\begin{figure}

{\centering \includegraphics{figures/soil-depth-plot-1} 

}

\caption{Mynd 3. Tegundatap sem fall af jarðvegsdýpi. Rauð lína = loess-sléttun; blá brotalína = línuleg aðlögun.}\label{fig:soil-depth-plot}
\end{figure}

\begin{center}\rule{0.5\linewidth}{0.5pt}\end{center}

\section{Jarðvegsbleyta --- kort yfir endurnýjun tegunda og
tegundatap}\label{jaruxf0vegsbleyta-kort-yfir-endurnuxfdjun-tegunda-og-tegundatap}

\begin{figure}

{\centering \includegraphics{figures/moisture-maps-1} 

}

\caption{Mynd 4. Endurnýjun tegunda (BC, efri kort) og breyting á tegundafjölda (ΔS, neðri kort) eftir jarðvegsbleytuflokkum. Rautt = mikil breyting / tap; grænt/blátt = stöðugleiki / aukning.}\label{fig:moisture-maps-1}
\end{figure}
\begin{figure}

{\centering \includegraphics{figures/moisture-maps-2} 

}

\caption{Mynd 4. Endurnýjun tegunda (BC, efri kort) og breyting á tegundafjölda (ΔS, neðri kort) eftir jarðvegsbleytuflokkum. Rautt = mikil breyting / tap; grænt/blátt = stöðugleiki / aukning.}\label{fig:moisture-maps-2}
\end{figure}
\begin{table}[!h]
\centering
\caption{\label{tab:moisture-stats}Endurnýjun tegunda og tegundatap eftir jarðvegsbleytuflokkum. Kruskal--Wallis: BC $\chi^2$=13.12 p=0.0044; $\Delta$S $\chi^2$=11.06 p=0.0114.}
\centering
\begin{tabular}[t]{lrrr}
\toprule
Jarðvegsbleyta & n reitir & Miðgildi BC & Miðgildi ΔS\\
\midrule
\cellcolor{gray!10}{Þurrlendi} & \cellcolor{gray!10}{111} & \cellcolor{gray!10}{0.678} & \cellcolor{gray!10}{-2}\\
Mýri & 77 & 0.609 & -4\\
\cellcolor{gray!10}{Deiglendi} & \cellcolor{gray!10}{70} & \cellcolor{gray!10}{0.526} & \cellcolor{gray!10}{-1}\\
Flói & 45 & 0.441 & -1\\
\bottomrule
\end{tabular}
\end{table}

Þurrlendi sýndi hæstu endurnýjun tegunda (BC miðgildi = 0.678) en fremur
lítið tegundatap (ΔS miðgildi = -2), sem bendir til þess að
tegundasamsetningin sé að breytast --- sumar tegundir hverfa en aðrar
koma í staðinn --- án þess að heildarfjöldi tegunda minnki verulega.
Mýrar sýndu hins vegar mest tegundatap (ΔS = -4) án þess að
Bray--Curtis-fjarlægðin væri samsvarandi há, sem bendir til þess að
tegundir séu að hverfa án þess að nýjar taki við. Þetta eru, í
grundvallaratriðum, tveir ólíkir breytingarferlar sem skipta máli fyrir
túlkun niðurstaðnanna. ---

\section{Vistgerðir --- kort yfir endurnýjun tegunda og
tegundatap}\label{vistgeruxf0ir-kort-yfir-endurnuxfdjun-tegunda-og-tegundatap}

\begin{figure}

{\centering \includegraphics{figures/habitat-maps-1} 

}

\caption{Mynd 5. Endurnýjun tegunda (BC, efri kort) og breyting á tegundafjölda (ΔS, neðri kort) eftir helstu vistgerðum (n ≥ 8 reitir).}\label{fig:habitat-maps-1}
\end{figure}
\begin{figure}

{\centering \includegraphics{figures/habitat-maps-2} 

}

\caption{Mynd 5. Endurnýjun tegunda (BC, efri kort) og breyting á tegundafjölda (ΔS, neðri kort) eftir helstu vistgerðum (n ≥ 8 reitir).}\label{fig:habitat-maps-2}
\end{figure}
\begin{table}[!h]
\centering
\caption{\label{tab:habitat-stats}Endurnýjun tegunda og tegundatap eftir vistgerð (n $\geq$ 8). Kruskal--Wallis: BC $\chi^2$=81.05 p $<$ 0,001; $\Delta$S $\chi^2$=53.60 p $<$ 0,001.}
\centering
\begin{tabular}[t]{lrrr}
\toprule
Vistgerð & n reitir & Miðgildi BC & Miðgildi ΔS\\
\midrule
\cellcolor{gray!10}{Lyngmóavist á hálendi} & \cellcolor{gray!10}{8} & \cellcolor{gray!10}{0.802} & \cellcolor{gray!10}{-8.0}\\
Rústamýravist & 31 & 0.798 & -7.0\\
\cellcolor{gray!10}{Rekjuvist} & \cellcolor{gray!10}{17} & \cellcolor{gray!10}{0.682} & \cellcolor{gray!10}{-4.0}\\
Starungsmýravist & 34 & 0.543 & -3.5\\
\cellcolor{gray!10}{Sjávarfitjungsvist} & \cellcolor{gray!10}{27} & \cellcolor{gray!10}{0.379} & \cellcolor{gray!10}{1.0}\\
\addlinespace
Mosahraunavist & 8 & 0.182 & 7.0\\
\bottomrule
\end{tabular}
\end{table}

Þrjár vistgerðir skera sig úr. \textbf{Rústamýravist} (n = 31) sýndi
bæði mikla endurnýjun tegunda (BC = 0.798) og mikið tegundatap (ΔS =
-7), og er þannig sú vistgerð sem er hvað verst sett samkvæmt báðum
mælikvörðum. \textbf{Lyngmóavist á hálendi} sýndi svipaðar niðurstöður
(BC = 0.802, ΔS = -8) og endurspeglar líklega hæðaráhrifin sem komu fram
í hlutfylgnigreiningu. Á hinn bóginn stendur \textbf{Mosahraunavist} upp
úr sem eina vistgerðin sem sýnir bæði lága endurnýjun tegunda (BC =
0.182) og aukningu á tegundafjölda (ΔS = 7).

\begin{center}\rule{0.5\linewidth}{0.5pt}\end{center}

\end{document}
